% Section: Budget Pacing Under Cost Drift
% Include from paper/sections/results.tex

\subsection{Budget Pacing Under Cost Drift (Figure~\ref{fig:exp03_adaptation})}
\label{sec:budget_drift}

\paragraph{Motivation.}
Experiments 01 and 02 studied budget pacing and non-stationarity in
isolation.  In production, both occur simultaneously: a router
operates under a cost budget when a model provider drops its pricing.
We ask: \emph{does the BudgetPacer automatically exploit a mid-stream
price drop to improve quality while maintaining budget compliance,
and how does it compare with static cost-control baselines?}

\paragraph{Setup.}
The $K{=}3$ portfolio follows a \emph{train-then-evaluate} pipeline
(matching Experiment~01).
The router first online-learns on the validation split
($n{=}1{,}785$) under normal pricing without recording evaluation
metrics; priors are trained on the disjoint training split
($n_{\mathrm{eff}}{=}10$, $\alpha{=}0.1$;
Section~\ref{sec:hparam_sweep}).
Evaluation uses the held-out test split ($n{=}1{,}824$),
partitioned into two phases:
Phase~1 ($n{=}912$) uses normal pricing (Gemini normalized cost
$\tilde{c}{=}0.67$); Phase~2 ($n{=}912$) applies a Gemini price
drop to \$0.10/\$0.10 per million tokens ($\tilde{c} \approx 0.0$).

Three budget targets are chosen to span qualitatively distinct
operating regimes of the primal-dual controller:
\begin{itemize}[nosep,leftmargin=*]
  \item \textbf{Tight} ($B{=}\$2.3{\times}10^{-4}$/req):
    The budget constraint fully dominates routing decisions.
    $\lambda_t$ is persistently high, forcing the router into a
    near-corner solution (Llama-dominant).  This regime tests
    whether the pacer can maintain exact budget compliance when
    there is almost no room to trade cost for quality.
  \item \textbf{Moderate} ($B{=}\$6.6{\times}10^{-4}$/req):
    The constraint is actively binding but not saturating.
    The router must continuously balance quality and cost,
    producing a mixed Llama/Mistral policy.  This is the
    regime where the closed-loop nature of the BudgetPacer
    matters most, because small routing changes cause immediate
    cost feedback that the dual update must absorb.
  \item \textbf{Loose} ($B{=}\$1.9{\times}10^{-3}$/req):
    The budget barely constrains routing; $\lambda_t$ is low
    and the router operates close to its unconstrained optimum.
    This regime tests whether the pacer gracefully ``gets out
    of the way'' when the budget is not a binding constraint.
\end{itemize}

\noindent
Together, these three levels cover the full spectrum from
budget-dominated to quality-dominated routing, allowing us to
evaluate adaptation behavior across the entire operating envelope
rather than at a single hand-picked budget.

\noindent
Per target, three conditions represent increasing levels of
routing sophistication, each using warmup priors:
\begin{enumerate}[nosep,leftmargin=*]
  \item \textbf{Fixed Policy (offline):}
    Frozen warmup priors with a matched static cost penalty;
    no online learning.  The industry standard: deploy an
    offline-trained router with a hand-tuned cost weight.
  \item \textbf{Naive Bandit ($\gamma{=}1.0$):}
    LinUCB with infinite memory and the same matched static
    cost penalty.  Online learning, but no principled budget
    mechanism and no forgetting.
  \item \textbf{BanditGPT (Pacer + $\gamma{=}0.997$):}
    Warmup priors with jointly-tuned forgetting and the
    primal-dual BudgetPacer targeting the dollar budget directly.
\end{enumerate}

\noindent
An \textbf{unconstrained BanditGPT} baseline (same algorithm,
same priors, no budget pacer, $\lambda_s{=}0$) is included in
every panel of Figure~\ref{fig:exp03_adaptation}.  This baseline
answers two questions that budget-constrained results alone cannot:
(i)~\emph{how much quality does the budget cost us?}---the gap
between a constrained curve and the unconstrained curve quantifies
the price of budget compliance; and
(ii)~\emph{does the pacer converge toward unconstrained behavior
after the price drop?}---when Gemini becomes nearly free, a
well-adapted pacer should approach the unconstrained routing mix,
and the unconstrained trace provides the reference to verify this.

\paragraph{Fixed Policy cost awareness.}
The ``Fixed Policy'' freezes the \emph{reward model} (no online
learning) but retains live cost awareness: the router recomputes
normalized costs from the model registry at each step.
When the Gemini price drop is applied, all conditions---including
the Fixed Policy---observe the new pricing, mirroring the real-world
scenario in which price changes are publicly announced.
The Fixed Policy can therefore shift routing in response to cost
signals (e.g., from Llama to Gemini after the drop) but
\emph{cannot} update its quality estimates for any arm.

\paragraph{Static cost-penalty calibration.}
The matched static cost penalties ($\lambda_s \in \{0.10, 0.30,
0.50\}$) were calibrated against the BudgetPacer's Phase~1 spend
profile.
Because the Fixed Policy and Naive Bandit route differently from
the pacer, the same penalty produces different actual spend;
this is a fundamental limitation of open-loop cost control and is
the central result of this experiment rather than a confound.
A grid search over $\lambda_s$ per condition would require running
each condition first---precisely the closed-loop calibration that
the BudgetPacer provides automatically.

% =====================================================================
% BUDGET COMPLIANCE (THE KEY DIFFERENTIATOR)
% =====================================================================

\paragraph{Budget compliance is the key differentiator.}
The defining advantage of the BudgetPacer is \emph{precise budget
targeting}:

\begin{center}
\small
\begin{tabular}{@{}llrr@{}}
\toprule
\textbf{Budget} & \textbf{Condition} & \textbf{Phase~1 Cost} & \textbf{Ratio} \\
\midrule
\multirow{3}{*}{Tight ($\$2.3 \times 10^{-4}$)}
  & Fixed Policy     & \$2.9e-5 & $0.12\times$ \\
  & Naive Bandit     & \$2.2e-4 & $0.92\times$ \\
  & \textbf{BanditGPT} & \textbf{\$2.4e-4} & $\mathbf{1.00\times}$ \\[3pt]
\multirow{3}{*}{Moderate ($\$6.6 \times 10^{-4}$)}
  & Fixed Policy     & \$1.3e-4 & $0.20\times$ \\
  & Naive Bandit     & \$7.0e-4 & $1.05\times$ \\
  & \textbf{BanditGPT} & \textbf{\$6.6e-4} & $\mathbf{1.00\times}$ \\[3pt]
\multirow{3}{*}{Loose ($\$1.9 \times 10^{-3}$)}
  & Fixed Policy     & \$9.2e-4 & $0.49\times$ \\
  & Naive Bandit     & \$5.2e-3 & $2.78\times$ \\
  & \textbf{BanditGPT} & \textbf{\$1.8e-3} & $\mathbf{0.97\times}$ \\
\bottomrule
\end{tabular}
\end{center}

\noindent
The \textbf{Ratio} column shows Phase~1 cost as a multiple of the
budget target.  Two failure modes emerge:

\begin{itemize}[nosep,leftmargin=*]
  \item \textbf{Fixed Policy under-spends} ($0.12$--$0.49\times$).
    The static cost penalty, tuned to produce similar Phase~1 spend
    for the pacer conditions, is too aggressive for a frozen policy
    that cannot refine its reward estimates.
    At tight budget, it routes 100\% to Llama-8B---the cheapest arm
    but also the lowest-quality---wasting the entire quality headroom
    that the budget allows.

  \item \textbf{Naive Bandit is inconsistent across budgets.}
    After pre-training on the validation split, the Naive Bandit
    happens to land near target at tight ($0.92\times$) and moderate
    ($1.05\times$) budgets---but this is coincidental, not controlled.
    At loose budget, the same static penalty ($\lambda_s{=}0.10$)
    cannot contain the shift toward expensive arms: the Naive Bandit
    spends \$5.2e-3/req ($2.78\times$ target).
    The static penalty provides no feedback mechanism to correct this;
    a different penalty per budget tier would require the very
    closed-loop calibration that the BudgetPacer provides automatically.
\end{itemize}

\noindent
\textbf{BanditGPT is the only condition that reliably hits the
budget target} ($0.97$--$1.00\times$).  The BudgetPacer's
primal-dual mechanism adjusts $\lambda_t$ continuously based on
observed costs, providing a closed-loop budget controller that
static penalties cannot replicate.

% =====================================================================
% ADAPTATION DYNAMICS
% =====================================================================

\paragraph{Why show the dual variable.}
The dual variable $\lambda_t$ is the central internal signal of the
BudgetPacer: it acts as a learned, time-varying ``price of cost''
that modulates the score of every arm at every step
(Eq.~\ref{eq:dual_update}).
When $\lambda_t$ is high, the router penalizes expensive arms
heavily; when $\lambda_t$ drops toward zero, cost is effectively
ignored and the router optimizes for quality alone.
Plotting $\lambda_t$ therefore exposes the \emph{mechanism} by
which the pacer adapts---not just the outcome (routing mix, cost)
but the control signal that produces it.
This is important because, unlike static cost penalties, $\lambda_t$
is never manually tuned: it emerges from the interaction between
the primal routing decisions and the dual budget feedback.
Observing its trajectory lets us verify that the adaptation is
smooth, timely, and regime-appropriate.

\paragraph{Lambda adaptation (Figure~\ref{fig:exp03_adaptation}a).}
The dual variable $\lambda_t$ responds directly to the cost EMA.
At \textbf{moderate} budget, $\lambda_t$ averages
$0.32$ during Phase~1 then drops sharply
to ${\sim}0.01$ within Phase~2 as
$\bar{c}_t$ falls below target.
At \textbf{loose} budget, $\lambda_t$ starts at $0.10$
and rapidly reaches near-zero after the price drop.
At \textbf{tight} budget, $\lambda_t$ averages $0.45$
in Phase~1 and stays elevated at $0.42$ in Phase~2---the budget
is so stringent that even the EMA takes many observations to
ratchet down from the near-ceiling values accumulated during
Phase~1.
This EMA-mediated response requires no external change-detection
signal: the dual update (Eq.~\ref{eq:dual_update}) naturally tracks
the cost regime.

\paragraph{Gemini adoption (Figure~\ref{fig:exp03_adaptation}b).}
At \textbf{moderate} budget, the pacer increases Gemini selection
from ${\sim}2\%$ to ${\sim}44\%$ within Phase~2, approaching the
unconstrained baseline (${\sim}75\%$ Gemini in Phase~2).
At \textbf{loose} budget, Gemini selection rises from ${\sim}9\%$
to ${\sim}48\%$, also converging toward the unconstrained level.
At \textbf{tight} budget, Gemini selection remains near $0\%$
even after the price drop---an important limitation discussed below.

\paragraph{Budget compliance (Figure~\ref{fig:exp03_adaptation}c).}

\begin{center}
\small
\begin{tabular}{@{}lrrrrr@{}}
\toprule
\textbf{Budget} & \textbf{Target} & \textbf{Phase~1} & \textbf{Phase~2} & \textbf{Overall} & \textbf{Util.} \\
\midrule
Tight    & \$2.34e-4 & \$2.35e-4 & \$2.34e-4 & \$2.34e-4 & $1.00\times$ \\
Moderate & \$6.62e-4 & \$6.59e-4 & \$3.67e-4 & \$5.13e-4 & $0.77\times$ \\
Loose    & \$1.87e-3 & \$1.82e-3 & \$3.53e-4 & \$1.09e-3 & $0.58\times$ \\
\midrule
Unconstr.\ & ---     & \$1.19e-2 & \$3.14e-4 & \$6.13e-3 & --- \\
\bottomrule
\end{tabular}
\end{center}

\noindent
BanditGPT achieves Phase~1 costs within 3\% of the target at all
budget levels.  After the price drop, costs decline below target
(the pacer constrains from above, never inflates); the overall
utilization ratio therefore falls below $1.0\times$.
The unconstrained baseline spends ${\sim}6$--$51\times$ more than
pacer targets in Phase~1 (\$1.19e-2/req), confirming that the pacer
provides genuine cost control.

\paragraph{Reward lift after price drop.}
We test whether the pacer exploits the price drop by measuring
$\Delta_\mathrm{reward} = \bar{r}_\mathrm{Phase\,2} -
\bar{r}_\mathrm{Phase\,1}$ within BanditGPT (a paired
comparison; each budget level serves as its own control).

\begin{center}
\small
\begin{tabular}{@{}lcccc@{}}
\toprule
\textbf{Budget} & \textbf{P1 Rwd} & \textbf{P2 Rwd} & $\boldsymbol{\Delta}$ & \textbf{95\% CI} \\
\midrule
Tight    & 0.854 & 0.861 & $+0.007$ & $[0.006,\; 0.009]$ \\
Moderate & 0.882 & 0.922 & $+0.041$ & $[0.038,\; 0.043]$ \\
Loose    & 0.907 & 0.924 & $+0.017$ & $[0.015,\; 0.019]$ \\
\bottomrule
\end{tabular}
\end{center}

\noindent
All budget levels show highly significant reward improvement
($p < 10^{-9}$ for all; one-sample $t$-test on per-seed deltas,
49~d.f.).  The moderate target exhibits the largest lift
($\Delta{=}+0.041$), reflecting the biggest shift from
Llama-dominated to Gemini-heavy routing.

\paragraph{Exploration starvation at tight budgets.}
At the tightest target, Gemini selection remains near 0\% even
after the price drop (Figure~\ref{fig:exp03_adaptation}b).
Two factors combine:
\begin{enumerate}[nosep,leftmargin=*]
  \item \textbf{EMA momentum.}
    $\lambda_t$ was saturated near $0.45$ for much of
    Phase~1.  The smoothing constant $\alpha_\mathrm{ema}{=}0.05$
    creates ${\sim}14$-step inertia; many low-cost observations are
    needed to ratchet $\lambda_t$ down from its Phase~1 plateau.
  \item \textbf{Information deficit.}
    With $<1\%$ Gemini selection during Phase~1, the bandit's reward
    model for Gemini is almost entirely the warmup prior---it has not
    learned enough online to build a confident, calibrated estimate.
    Even when the cost penalty vanishes in Phase~2, the LinUCB
    exploration bonus is too small (with $n_\mathrm{eff}{=}10$
    prior) to pull Gemini above the well-explored Llama and Mistral
    arms.
\end{enumerate}

\noindent
This is a known property of budget-constrained bandits: very tight
budgets suppress exploration of expensive arms, and the resulting
information deficit persists even after the cost barrier is removed.

\paragraph{Synthesis.}
This experiment demonstrates that
\textbf{static cost penalties cannot substitute for principled
budget pacing}:
\begin{itemize}[nosep,leftmargin=*]
  \item The \textbf{Fixed Policy} under-spends the budget by
    $2$--$8\times$ because it cannot refine its reward estimates---it
    sacrifices quality unnecessarily.
  \item The \textbf{Naive Bandit} is inconsistent: it coincidentally
    tracks the target at tight and moderate budgets but over-spends
    by $2.8\times$ at loose budget, because no feedback mechanism
    corrects the static cost penalty.
  \item \textbf{BanditGPT} is the only condition that achieves
    precise budget compliance ($0.97$--$1.00\times$ target) while
    automatically exploiting the Gemini price drop to improve quality
    ($\Delta_\mathrm{reward}{=}+0.007$ to $+0.041$).
  \item Very tight budgets create an \textbf{exploration starvation}
    effect: the bandit cannot learn about expensive arms it was
    prevented from selecting, and this information deficit persists
    after the cost barrier is removed.
\end{itemize}
